\chapter{Introduction}\label{chap:introduction}
This work was motivated by the need for anatomically accurate, patient-specific heart valve phantoms for pre-procedural planning and device testing. The following sections outline the motivation, aim, research question, and scientific contribution of this thesis. The data which this work heavily relies on stems from Siemens Healthcare Diagnostics GmbH (\emph{Siemens Healthineers}), who provided the volumetric ultrasound data and the unthickened surface mesh of the mitral valve leaflets. The pseudonymised imaging data was acquired from a patient with mitral regurgitation. The malfunction in the valve was deliberately kept within the data to ensure that the resulting phantom is clinically relevant and can be used i.e. for flow simulations of a patient specific diseased mitral valve. 

Besides the ultrasound data of the valves Siemens Healthineers also provided the imaging data of the atria and ventricles of the heart from the same patient. However, since the substance of producing a physical phantom prototype of the mitral valve leaflets is already complex enough, this work focuses on the leaflets only. A similar task of producing a physical phantom of the atria and ventricles is done by a colleague and fellow student \emph{Peter Machoritsch} at the Technical University of Vienna in parallel to this work. The two works are independent of each other, however some parts especially in Section~\ref{chap:methodology} may overlap. The resulting physical phantoms should be combinable to form a complete heart phantom which is the goal of the overall project. It should be emphasized that whenever certain design-choices with respect to eachothers works are made, they are documented and justified to ensure reproducibility and scientific rigor. 


\section{Motivation}\label{sec:intro_motivation}

Heart valve diseases, particularly mitral regurgitation and aortic stenosis, represent a major global health burden. Echocardiography - both transthoracic (TTE) and transesophageal (TEE) - remains the clinical gold standard for non-invasive evaluation of valvular anatomy, leaflet motion, and hemodynamic severity. However, ultrasound imaging is inherently operator-dependent and subject to intra- and inter-observer variability, image artifacts, and limited acoustic windows.

To improve diagnostic accuracy, train clinicians, and benchmark automated quantification algorithms, physical cardiac phantoms are increasingly utilized. Traditionally, these phantoms serve as standardized test objects for flow measurement and leaflet segmentation. More recently, a critical second application has emerged: patient-specific device implantation planning. For complex structural heart interventions, such as Transcatheter Mitral Valve Replacement (TMVR) or transcatheter edge-to-edge repair (TEER), anatomically accurate physical phantoms enable pre-procedural fit assessment, mechanical tissue-device interaction testing, and post-implantation flow simulation.

In this context, the geometric accuracy of the phantom - specifically the spatial distribution of leaflet thickness - plays a crucial role. Heart valve leaflets are not uniform surfaces; their thickness varies continuously across different anatomical zones (e.g., from the thin leaflet center to the thicker outer rim). Approximating leaflets with a uniform wall thickness compromises both the mechanical compliance of the model and its acoustic reflection characteristics under ultrasound. Creating an anatomically precise, thickness-accurate phantom directly from clinical volumetric ultrasound data remains a technical challenge due to speckle noise, low contrast-to-noise ratios, and spatial resolution constraints.

\section{Aim of the Work}\label{sec:intro_aim}

The primary aim of this work is to design and produce a 3D-printed prototype of an anatomically accurate mitral valve ultrasound phantom with realistic, spatially varying leaflet thickness. Starting from a volumetric ultrasound scan and an unthickened single-surface leaflet mesh, this project establishes an end-to-end processing pipeline to recover local leaflet thickness from image data, transfer it onto the geometric mesh, construct a closed solid model, and manufacture a flexible physical phantom suitable for ultrasound imaging and device interaction assessment.

\section{Research Question}\label{sec:intro_rq}

To guide the development and evaluation of the pipeline, this thesis addresses the following central research question:

\begin{quote}
\emph{How can spatially varying leaflet thickness be recovered from volumetric ultrasound data and transferred to a 3D-printable, anatomically accurate valve model?}
\end{quote}

\section{Scientific Contribution}\label{sec:intro_contribution}

The key contributions of this work are:

\begin{enumerate}
    \item \textbf{Hand-Guided, Locally Validated Thickness Recovery:} A hybrid measurement strategy that avoids the pitfalls of global volumetric thresholding in noisy ultrasound data by combining local intensity profile criteria with guided manual sampling in reliable image regions, supplementing missing values by robust median estimates in low-quality regions.
    \item \textbf{Multi-Software Bounding-Box Alignment:} A reproducible geometric registration method using bounding-box reference planes to align disparate coordinate spaces across MeVisLab, 3D Slicer, and Blender without relying on ambiguous curved and complex surface landmarks.
    \item \textbf{Spatially Interpolated Boundary Surface Extrusion:} A method to map discrete thickness measurements across complex surface meshes using inverse distance weighting (IDW) and single-sided boundary surface extrusion, preserving authentic anatomical thickness gradients without artificial uniform wall assumptions.
\end{enumerate}